\chapter{Implementación de la solución}
\label{ch:implementacion}
\section{Componentes y tecnologías}
\Analyzer{} es una herramienta en Python que, después del workflow, combina el directorio de Pegasus con el historial exportado de HTCondor. Produce resúmenes y desgloses de obligaciones, entregas, evidencia y ubicación; cada análisis conserva la versión y la identidad del flujo.

El cálculo incorpora el relevo del capítulo~\ref{ch:diseno}: si un intermedio tiene consumidor en otro worker, añade la obligación al coordinador y la deduplica cuando este también es destino final. Si no se identifica el coordinador, marca incompleta la evidencia.

\paragraph{Integración de las fuentes.}
El grafo expresa tareas y archivos; el plan añade actividades de preparación y entrega y puede reunir tareas en un job. La lista secuencial, los comandos y las dependencias permiten reconstruir cada grupo; el historial asigna a sus tareas la máquina del job. La herramienta coteja flujo e historial, identifica el intento y comprueba terminación y eventos según el perfil.

Catálogos y metadatos permiten conciliar tamaños e identidades. Fuentes incompatibles, colisiones de rutas, ubicación ausente o tareas repetidas impiden reconstruir la parte afectada; nombres similares no bastan para establecer identidad.

\section{Conciliación de entregas}
Una dirección corresponde a la entrada de archivos al job o a su salida. Para cada job $j$ y dirección $d$, se reconstruye la lista efectiva de archivos científicos $S_{jd}$ y se separan los auxiliares. Si se conocen sus tamaños, el residuo de control se calcula restando del total nativo $b_{jd}$ los bytes auxiliares $a_{jd}$ y los tamaños de los archivos científicos:
\begin{equation}
r_{jd}=b_{jd}-a_{jd}-\sum_{f\in S_{jd}}s(f).
\label{eq:residuo}
\end{equation}
El perfil del bloque A exige residuo cero, conteos concordantes y consistencia con terminación e intento. Se contrastan los totales nativos de HTCondor por dirección con el último intento y se corrobora la lista mediante evidencia agregada del job. El caso GL3 del capítulo~\ref{ch:validacion} muestra la cuenta completa.

Si se desconoce algún tamaño auxiliar, el modo \texttt{SUCCESSFUL\_MANIFEST} exige manifiesto compatible, tamaños científicos y ejecución satisfactoria; el residuo queda indeterminado. La auditoría histórica del bloque A usa \texttt{EXACT\_BYTES}. Los archivos sin evidencia suficiente se conservan como faltantes para decidir el cálculo.

\paragraph{Reportes y control de consistencia.}
El resumen presenta $M_{\mathrm{req}}$, $M_{\mathrm{obs}}$, DME y estado; los desgloses identifican destinos, entregas por dirección y conciliación de auxiliares. Los diagnósticos separan evidencia faltante, contradicciones y perfiles fuera de alcance.

La comparación con código separado comprueba consistencia, aunque comparta definición y fuentes.

\paragraph{Trazabilidad de una ejecución agrupada.}
\label{sec:traza_impl}
El grupo se reconstruye antes de calcular ambos volúmenes. En GL3, plan e historial sitúan tres tareas en un job de W1, y el grafo conserva dos dependencias intermedias. La rama de obligaciones asigna origen y destino en W1 a esos archivos; la de entregas recoge los intermedios y el final declarados en el contrato de salida. Así, un intermedio puede no crear una obligación entre máquinas y aun aportar al observado mediante un retorno corroborado.



El siguiente capítulo verifica estas cuentas y examina qué entregas cambian al modificar agrupación o contratos.
